\section{Mukai--Gram dictionary on \texorpdfstring{$K3\times E$}{K3 x E}}
\label{appendix:mukai-gram-dictionary}

Let \(X=S\times E\) be the smooth trivial-canonical threefold, with
\(S\) a smooth complex K3 surface and \(E\) a smooth complex elliptic curve. The
Mukai vector and Mukai pairing on \(\widetilde H(S,\Z)\), together with
the K\"unneth decomposition on \(H^*(X,\Z)\), produce the Gram-degree
projection \(\Pi_X:\Gamma_X^{\mathrm{form}}\to\Gamma_{\mathrm{gram}}\)
and its normal-ordered companion \(\overline\Pi_X\) on the central
extension \(\widehat\Gamma_X\). The Borcherds expansion variables
\(q^nr^ls^m\) of an Igusa cusp expansion record Gram-degree shadows in
\(\Gamma_{\mathrm{gram}}\).  When the shadow is carried by an actual
one-dimensional subscheme \(Y\subset X\), its dyonic charge
\(c=(Q_Y,P_Y)\) satisfies the formula below.  The existence of such a
subscheme is the algebraic-effective Hall-support condition, not a
formal consequence of the Fourier monomial. For \(Y\) of class
\([Y]=(\beta_h,d)\) and \(\chi(\mathcal O_Y)=n\),
\[
\Pi_X(Q_Y,P_Y)=(h-1,\,n,\,d-1).
\]
The threefold arithmetic Euler characteristic
\(n=\chi(\mathcal O_Y)\) is the direct structure-sheaf invariant; no
``Todd correction'' enters, and no ambient fourfold appears.

The Mukai--Gram construction sits at Beilinson chart layer
\(\alpha\) (intrinsic charge data) and
\(\gamma\) (ambient Mukai/Gram lattice). The
\(\kappa\)-tuple, orientation character \(\epsilon_o\), Pfaffian
section, and recognition theorem live at higher Beilinson levels and
depend on this chart datum.

\subsection{Mukai vector and Mukai pairing on a K3 surface}
\label{subsec:mukai-vector-pairing}

\begin{definition}[Mukai lattice and Mukai vector]
\label{def:mukai-lattice-vector-appendixC}
\ClaimStatusDefinition\
Let \(S\) be a smooth complex K3 surface. The Mukai lattice is the
total cohomology
\[
\widetilde H(S,\Z)
:=H^0(S,\Z)\oplus H^2(S,\Z)\oplus H^4(S,\Z),
\]
graded by even cohomological degree, an even unimodular lattice of
signature \((4,20)\) and rank \(24\) when equipped with the Mukai
pairing below. For a coherent sheaf \(F\in\mathrm{Coh}(S)\),
respectively for an object \(F\in D^b(S)\), the Mukai vector is
\[
v_S(F)
:=\ch(F)\cdot\sqrt{\td(S)}
=\bigl(\rk F,\;c_1(F),\;\ch_2(F)+\rk F\cdot[\mathrm{pt}]\bigr)
\;\in\;\widetilde H(S,\Z).
\]
The square root \(\sqrt{\td(S)}=(1,0,1)\) on the K3 surface uses the
Todd class \(\td(S)=(1,0,2[\mathrm{pt}])\), so
\(\sqrt{\td(S)}=(1,0,[\mathrm{pt}])\); identifying
\([\mathrm{pt}]\) with the integer \(1\) in the displayed coordinates
gives the displayed vector. The triple
\(v_S(F)=(r,D,s)\) records the rank \(r\in H^0\), the first Chern class
\(D\in H^2\), and the Mukai-corrected second component
\(s=\ch_2(F)+r\in H^4\). Mukai introduced this lattice and pairing in
\cite[Theorem~1.5]{Mukai1987}, with the symplectic structure of the
moduli space of sheaves established in \cite{Mukai1984}.
\end{definition}

\begin{proposition}[Mukai pairing on the K3 surface]
\label{prop:mukai-pairing-appendixC}
\ClaimStatusProvedElsewhere\
The Mukai pairing on \(\widetilde H(S,\Z)\) is the symmetric \(\Z\)-bilinear
form
\[
\bigl\langle(r,D,s),(r',D',s')\bigr\rangle_{\mathrm{Muk}}
=D\cdot D'-rs'-r's,
\]
equivalently
\[
\bigl\langle v,w\bigr\rangle_{\mathrm{Muk}}
=-\int_S v^{\vee}\cup w,
\qquad
v^\vee:=(r,-D,s),
\]
where \(v_2\cup w_2\) is the K3 cup product on \(H^2(S,\Z)\) and
\(v_0\cup w_4\), \(v_4\cup w_0\) record the cross-degree pairings.
The form is symmetric, even, unimodular, of signature \((4,20)\); its
hyperbolic-plane summands are
\[
\widetilde H(S,\Z)
\simeq U^{\oplus4}\oplus E_8(-1)^{\oplus2}
\]
as integral lattices.  For a polarized K3 of Picard rank \(\rho\), the
N\'eron--Severi lattice \(\mathrm{NS}(S)\) has signature
\((1,\rho-1)\), while the algebraic Mukai lattice
\[
\widetilde H^{1,1}_{\mathrm{alg}}(S,\Z)
=H^0(S,\Z)\oplus\mathrm{NS}(S)\oplus H^4(S,\Z)
\]
has signature \((2,\rho)\).  This is Mukai
\cite[Theorem~1.5(ii)]{Mukai1987}; see also Huybrechts--Lehn
\cite[Chapter~6]{HuybrechtsLehn}.
\end{proposition}

\begin{proof}
The bilinear identity is the definition; its symmetry is immediate from
the symmetry of the K3 cup product and the identification of degree-zero
and degree-four components. Even unimodularity follows from the K3
intersection lattice on \(H^2\) (even unimodular of signature
\((3,19)\)) and the unimodular pairing
\(H^0\times H^4\to\Z\) under \(rs'\mapsto rs'\). The signature
\((4,20)\) follows by Hodge index plus the hyperbolic factor on
\(H^0\oplus H^4\).  The same hyperbolic factor changes
\(\mathrm{NS}(S)\) of signature \((1,\rho-1)\) into
\(\widetilde H^{1,1}_{\mathrm{alg}}(S,\Z)\) of signature
\((2,\rho)\).  The lattice isomorphism is
\cite[Theorem~1.5(ii)]{Mukai1987}.
\end{proof}

\begin{remark}[Mukai pairing is symmetric on both parts]
\label{rem:mukai-symmetric-both-parts}
\ClaimStatusProved\
The Mukai pairing on \(\widetilde H(S,\Z)\) is symmetric in both the
algebraic and transcendental parts. There is no orthosymplectic
\(\mathfrak{osp}\) refinement at the level of the lattice. The
Hodge \(\Z/2\)-super-extension \((\mathbb C,\omega_S)\mapsto
\mathbb C\oplus\omega_S\) used in Mukai self-mirror
language imposes a separate \(\Z/2\)-grading on a Yangian-type quantum
group constructed from this lattice; that grading does not change the
underlying lattice pairing, which remains symmetric. The classical Lie
limit of the Mukai-form Yangian is the orthogonal Lie algebra
\(\mathfrak{so}(4,20)\), not \(\mathfrak{osp}(4\mid20)\); see
\cite[\S2]{SchiffmannVasserot2013} for the Mukai-side Yangian
classical limit and \cite{KSCoHA} for the cohomological Hall algebra
realization.
\end{remark}

\begin{example}[Mukai vector of a point and of an ideal sheaf]
\label{ex:mukai-vector-point-and-ideal}
\ClaimStatusDerivation\
\(v_S(\mathcal O_p)=(0,0,1)\) for the structure sheaf of a point
\(p\in S\); \(v_S(\mathcal O_S)=(1,0,1)\) for the trivial line bundle;
\(v_S(L)=(1,c_1(L),c_1(L)^2/2+1)\) for a line bundle \(L\). For an ideal
sheaf \(\mathcal I_C\) of a smooth curve \(C\subset S\) of genus \(g\),
the adjunction formula on the K3 surface gives
\(C^2=c_1(\mathcal O_S(C))^2=2g-2\), and
\[
v_S(\mathcal I_C)=(1,-c_1(\mathcal O_S(C)),\ch_2(\mathcal I_C)+1).
\]
The Mukai pairing
\(\langle v_S(\mathcal O_p),v_S(\mathcal O_p)\rangle_{\mathrm{Muk}}
=0-0\cdot1-0\cdot1=0\), recording that points have zero self-pairing in
the Mukai form. The pairing
\(\langle v_S(\mathcal O_S),v_S(\mathcal O_p)\rangle_{\mathrm{Muk}}
=0-1\cdot1-0\cdot0=-1\) records the Euler characteristic
\(\chi(\mathcal O_S,\mathcal O_p)=1\) up to the Mukai sign convention
\(\chi(F,G)=-\langle v_S(F),v_S(G)\rangle_{\mathrm{Muk}}\) of
\cite[Proposition~2.2]{Mukai1987}.
\end{example}

\subsection{K\"unneth on \texorpdfstring{$K3\times E$}{K3 x E} and the rank-3 Gram lattice}
\label{subsec:kunneth-k3-times-e-appendixC}

\begin{proposition}[K\"unneth for \(X=K3\times E\)]
\label{prop:kunneth-k3-times-e-appendixC}
\ClaimStatusProvedElsewhere\
Let \(S\) be a smooth complex K3 surface and \(E\) a smooth complex
elliptic curve. The even integral cohomology of \(X=S\times E\) splits as
\[
H^{\mathrm{even}}(X,\Z)
\simeq H^*(S,\Z)\otimes_{\Z}H^{\mathrm{even}}(E,\Z),
\]
the odd--odd K\"unneth term \(H^{\mathrm{odd}}(S)\otimes H^{\mathrm{odd}}(E)\)
absent because \(S\) carries no odd cohomology, both factors free
\(\Z\)-modules. Fixing a generator \(\omega_E\in H^2(E,\Z)\) with
\(\int_E\omega_E=1\),
\(H^{\mathrm{even}}(E,\Z)=\Z\cdot1_E\oplus\Z\cdot\omega_E\), and every
even class on \(X\) decomposes uniquely as
\[
v_X=v^{(0)}\otimes 1_E+v^{(\omega)}\otimes\omega_E,
\qquad v^{(0)},v^{(\omega)}\in H^*(S,\Z).
\]
The formal K\"unneth lattice
\[
H^{\mathrm{even}}(X,\Z)
\simeq \widetilde H(S,\Z)\otimes U_E,\qquad
U_E:=H^0(E,\Z)\oplus H^2(E,\Z),
\]
carries the tensor K\"unneth pairing of rank \(48\) and signature
\((24,24)\), obtained from the \((4,20)\)-Mukai pairing on
\(\widetilde H(S,\Z)\) and the \((1,1)\)-hyperbolic pairing on \(U_E\).
This cross-pairing is distinct from the direct-sum componentwise Gram
pairing on \(\Gamma_X^{\mathrm{form}}\) used by \(\Pi_X\).  The odd
K\"unneth term \(H^{\mathrm{odd}}(S)\otimes H^{\mathrm{odd}}(E)\) is
absent because the K3 surface \(S\) carries no odd cohomology. See
Huybrechts \cite[Chapter~3]{HuybrechtsLehn}.
\end{proposition}

\begin{proof}
\(E\) and \(S\) have torsion-free integral cohomology, so K\"unneth holds
integrally and produces the displayed tensor decomposition. The
tensor-pairing factorization is bilinearity of the cup product over the
K\"unneth isomorphism. The lattice signature comes from
\((4,20)\) on \(\widetilde H(S,\Z)\) (Proposition~\ref{prop:mukai-pairing-appendixC})
tensored with \((1,1)\) on \(U_E\).
\end{proof}

\begin{definition}[Gram-degree projection on \(X=K3\times E\)]
\label{def:gram-degree-projection-appendixC}
\ClaimStatusDefinition\
The formal additive even-Mukai sector is
\[
\Gamma_X^{\mathrm{form}}
:=\widetilde H(S,\Z)\oplus\widetilde H(S,\Z),\qquad
c=(Q,P)\in\Gamma_X^{\mathrm{form}},
\]
with \(Q,P\in\widetilde H(S,\Z)\). The Gram-index lattice is
\[
\Gamma_{\mathrm{gram}}:=\Z^3,
\]
and the Gram-degree projection is the quadratic map
\[
\Pi_X:\Gamma_X^{\mathrm{form}}\longrightarrow\Gamma_{\mathrm{gram}},
\qquad
(Q,P)\longmapsto
\Bigl(\tfrac{(Q,Q)_{\mathrm{Muk}}}2,\,
(Q,P)_{\mathrm{Muk}},\,
\tfrac{(P,P)_{\mathrm{Muk}}}2\Bigr).
\]
The triple \((n,l,m)=\Pi_X(Q,P)\) is the Gram-degree shadow; it is the
target of the Igusa Fourier expansion at the standard cusp. It is
\emph{not} an additive Hall grading: the additivity defect is the
symmetric bilinear cocycle \(B\) of \S\ref{subsec:gram-cocycle-appendixC}
below.
\end{definition}

\subsection{D6--D2--D0 dictionary}
\label{subsec:d6-d2-d0-dictionary-appendixC}

The K\"unneth decomposition of \(v_X(\mathcal I_Y)\) on the threefold
\(X=K3\times E\) produces the D6--D2--D0 dictionary
\(\Pi_X(Q_Y,P_Y)=(h-1,n,d-1)\) of
Theorem~\ref{thm:d6-d2-d0-dictionary-appendixC} as a chart-data
identity. The Mukai-vector calculation isolates the K3-fibre rank-one
component \(P_Y=(1,0,1-d)\) on the \(1_E\)-summand and the curve-class
component \(Q_Y=(0,-\beta^\vee,-n)\) on the \(\omega_E\)-summand, with
\(\beta^\vee\in H^2(S,\Z)\) the Poincar\'e dual of \(\beta\); the Mukai
pairing on the K3 components delivers the Gram-degree shadow.

\begin{theorem}[Mukai vector of an ideal sheaf on the threefold $K3\times E$]
\label{thm:mukai-vector-ideal-sheaf-threefold}
\ClaimStatusProvedHere\
Let \(X=S\times E\) with \(S\) a smooth complex K3 surface and \(E\) a
smooth complex elliptic curve, so \(K_X\simeq\mathcal O_X\) and \(X\)
is a smooth trivial-canonical threefold. Let \(Y\subset X\) be a
one-dimensional closed subscheme of
class \([Y]=(\beta,d)\in H_2(S,\Z)\oplus H_2(E,\Z)\), write
\(\beta^\vee\in H^2(S,\Z)\) for the Poincar\'e dual of \(\beta\), and let
\(n=\chi(\mathcal O_Y)\in\Z\). The Mukai vector of the
ideal sheaf \(\mathcal I_Y\in D^b(X)\) decomposes against the K\"unneth
basis as
\[
v_X(\mathcal I_Y)
=(1,0,1-d)\otimes1_E
+(0,-\beta^\vee,-n)\otimes\omega_E.
\]
\end{theorem}

\begin{proof}
The trivial-canonical threefold condition gives
\(\sqrt{\td(X)}=\pi_S^*\sqrt{\td(S)}\,\pi_E^*\sqrt{\td(E)}=(1,0,1)\otimes1\),
since the elliptic curve is one-dimensional with trivial Todd class
\(\td(E)=1\). The ideal-sheaf identity is
\(\ch(\mathcal I_Y)=1-\ch(\mathcal O_Y)\). For the structure sheaf of a
one-dimensional subscheme of a trivial-canonical threefold,
\(\ch_0(\mathcal O_Y)=0\), \(\ch_1(\mathcal O_Y)=0\),
\(\ch_2(\mathcal O_Y)=\operatorname{PD}[Y]\), and
\(\int_X\ch_3(\mathcal O_Y)=\chi(\mathcal O_Y)=n\). Decomposing
\([Y]=(\beta,d)\) on the K\"unneth basis,
\(\operatorname{PD}[Y]=\beta^\vee\otimes\omega_E+d[\mathrm{pt}_S]\otimes 1_E\),
where \([\mathrm{pt}_S]\) is the K3 point class; the fibrewise
elliptic-degree component \(d[\mathrm{pt}_S]\otimes 1_E\) records the
elliptic degree. Hence
\[
\ch(\mathcal O_Y)
=(0,0,d)\otimes1_E+(0,\beta^\vee,n)\otimes\omega_E,
\]
and
\(1-\ch(\mathcal O_Y)
=(1,0,-d)\otimes1_E+(0,-\beta^\vee,-n)\otimes\omega_E\).
Multiplying by
\(\sqrt{\td(X)}=(1,0,1)\otimes1_E\) using the K3 cup product on each
K\"unneth summand,
\((1,0,1)\cup(1,0,-d)=(1,0,1-d)\) since \([\mathrm{pt}_S]\cup1=[\mathrm{pt}_S]\),
and \((1,0,1)\cup(0,-\beta^\vee,-n)=(0,-\beta^\vee,-n)\) since
\(\beta^\vee\cup[\mathrm{pt}_S]=0\) in \(H^*(S,\Z)\), yields the displayed
Mukai vector
\(v_X(\mathcal I_Y)=(1,0,1-d)\otimes1_E+(0,-\beta^\vee,-n)\otimes\omega_E\).
\end{proof}

\begin{remark}[Threefold Euler characteristic, no Todd correction]
\label{rem:threefold-euler-characteristic-no-todd-correction}
\ClaimStatusDerivation\
The integer \(n=\chi(\mathcal O_Y)\in\Z\) is the holomorphic Euler
characteristic of the structure sheaf \(\mathcal O_Y\) on the threefold
\(X=S\times E\); the trivial-canonical threefold condition makes
\(\sqrt{\td(X)}\) act as identity on the third Chern character so that
\(\int_X\ch_3(\mathcal O_Y)=n\) directly. There is no auxiliary
``Todd correction'' to install on top of \(n\); the Mukai vector
\((0,-\beta^\vee,-n)\otimes\omega_E\) records \(n\) in the K3
\(H^4\)-coordinate on the elliptic \(\omega_E\)-summand. The sign is fixed by
\(\ch(\mathcal I_Y)=1-\ch(\mathcal O_Y)\) and the K\"unneth orientation
\(\int_E\omega_E=+1\).
\end{remark}

\begin{theorem}[D6--D2--D0 Mukai--Gram dictionary, restated as chart data]
\label{thm:d6-d2-d0-dictionary-appendixC}
\ClaimStatusProvedHere\
With the standing data of
Theorem~\ref{thm:mukai-vector-ideal-sheaf-threefold}, set
\[
P_Y:=(1,0,1-d),\qquad
Q_Y:=(0,-\beta^\vee,-n)\;\in\;\widetilde H(S,\Z),
\]
the K\"unneth components of \(v_X(\mathcal I_Y)\). Then the Gram-degree
projection of Definition~\ref{def:gram-degree-projection-appendixC} reads
\[
\Pi_X(Q_Y,P_Y)
=\Bigl(\tfrac{(Q_Y,Q_Y)_{\mathrm{Muk}}}2,\,
(Q_Y,P_Y)_{\mathrm{Muk}},\,
\tfrac{(P_Y,P_Y)_{\mathrm{Muk}}}2\Bigr)
=\Bigl(\tfrac{(\beta^\vee)^2}2,\,n,\,d-1\Bigr).
\]
For the elliptically fibered K3 surface \(S\to\mathbb P^1\) with section
\(s_1\) and curve class
\(\beta=\beta_h=s_1+hF\) (\(h\in\Z_{\ge0}\), \(F\) the fibre class), the
self-intersection
\(\beta_h^2=2h-2\) gives
\[
\Pi_X(Q_Y,P_Y)=(h-1,\,n,\,d-1).
\]
\end{theorem}

\begin{proof}
Apply the Mukai pairing of
Proposition~\ref{prop:mukai-pairing-appendixC} component by component:
\begin{align*}
(Q_Y,Q_Y)_{\mathrm{Muk}}
&=(-\beta^\vee)\cdot(-\beta^\vee)-0\cdot(-n)-0\cdot(-n)=(\beta^\vee)^2,\\
(Q_Y,P_Y)_{\mathrm{Muk}}
&=(-\beta^\vee)\cdot 0-0\cdot(1-d)-1\cdot(-n)=n,\\
(P_Y,P_Y)_{\mathrm{Muk}}
&=0\cdot0-1\cdot(1-d)-1\cdot(1-d)=2(d-1),
\end{align*}
giving the half-pairings \(((\beta^\vee)^2/2,n,d-1)\). For the elliptic class
\(\beta_h=s_1+hF\), identified with its Poincar\'e dual, the K3 intersection \(s_1^2=-2\) on a section of an
elliptic K3 surface (genus formula \(2g(s_1)-2=s_1^2+s_1\cdot K_S=-2+0\),
giving \(g(s_1)=0\)), \(s_1\cdot F=1\), \(F^2=0\), so
\(\beta_h^2=s_1^2+2h(s_1\cdot F)+h^2 F^2=-2+2h=2h-2\). The Gram-degree
shadow follows. The cross-level identity with the OP/DT scalar
normalization is the
\(Q_Y^2/2=h-1\) line of Theorem~\ref{thm:d6-d2-d0-dictionary-appendixC}.
\end{proof}

\begin{remark}[Algebraic vs primitive D2 charges]
\label{rem:algebraic-versus-primitive-d2-charges}
\ClaimStatusDerivation\
The K3 curve class \(\beta\in H_2(S,\Z)\) appearing through its
Poincar\'e dual in \(Q_Y=(0,-\beta^\vee,-n)\) is algebraic:
\(\beta^\vee\in\mathrm{NS}(S)\), since \(Y\subset X\) is a
closed subscheme. Primitivity of \(\beta\) — that \(\beta\) is not a
non-trivial integer multiple of another effective class — is a separate
condition; the dictionary holds for any algebraic \(\beta\), not only
primitive ones. The Mukai-pairing formula
\((Q_Y,Q_Y)_{\mathrm{Muk}}=(\beta^\vee)^2\) records the
\(\mathrm{NS}(S)\)-self-intersection regardless of primitivity. The
Bryan--Oberdieck--Pandharipande--Yin/Maulik--Pandharipande conventions
\cite{OPand,BRYAN,MaulikTodaGV} take \(\beta\) algebraic; the elliptic
degree \(d\in\Z_{\ge0}\) is unconstrained beyond non-negativity in the
effective semigroup.
\end{remark}

\begin{remark}[Chart-data layer; not the additive Hall grading]
\label{rem:chart-data-not-additive-hall-grading}
\ClaimStatusProved\
The triple \((n,l,m)=\Pi_X(Q_Y,P_Y)\) records the Gram-degree
invariants of the dyonic charge \(c=(Q,P)\). It is the target of the
Igusa Fourier expansion and the Borcherds product variables
\(q^nr^ls^m\). It is \emph{not} the additive Hall grading on a
hypothetical compact \(K3\times E\) Hall sector; the additive grading is
the upstairs charge \(c\in\Gamma_X^{\mathrm{form}}\) (or its
normal-ordered lift \(\widehat\Gamma_X\) below), and \(\Pi_X\) is a
quadratic map onto Gram invariants. The polarization defect that obstructs
the upstairs additive grading from descending directly is the symmetric
bilinear cocycle \(B\) of \S\ref{subsec:gram-cocycle-appendixC}.
\end{remark}

\subsection{Symmetric bilinear Gram cocycle}
\label{subsec:gram-cocycle-appendixC}

\begin{lemma}[Gram additivity defect, restated as chart structure]
\label{lem:gram-additivity-defect-appendixC}
\ClaimStatusProvedElsewhere\
For \(c_i=(Q_i,P_i)\in\Gamma_X^{\mathrm{form}}\), \(i=1,2\),
\[
\Pi_X(c_1+c_2)
=\Pi_X(c_1)+\Pi_X(c_2)+B(c_1,c_2),
\]
where the symmetric bilinear cocycle \(B\) is
\[
B(c,c')
:=\bigl((Q,Q')_{\mathrm{Muk}},\;
(Q,P')_{\mathrm{Muk}}+(Q',P)_{\mathrm{Muk}},\;
(P,P')_{\mathrm{Muk}}\bigr).
\]
The map \(B\) is symmetric and \(\Z\)-bilinear with values in
\(\Gamma_{\mathrm{gram}}=\Z^3\), satisfies the polarization identity
\[
B(c,c)=2\Pi_X(c),
\]
and equals \(-\delta\Pi_X\) as an additive coboundary,
\((\delta\Pi_X)(c,c'):=\Pi_X(c)+\Pi_X(c')-\Pi_X(c+c')\). Thus the
ordinary additive cohomology class \([B]=0\); the obstruction recorded
here is the impossibility of an \emph{additive linear} trivialization,
not a non-zero additive cohomology class.
\end{lemma}

\begin{proof}
This is Lemma~\ref{lem:gram-additivity-defect}.  Symmetry,
bilinearity, integrality, polarization, and \(\delta\)-cocycle property
are Lemma~\ref{lem:gram-additivity-defect-appendixC}(i)--(iii). The
splitting \(B=-\delta\Pi_X\) by the quadratic refinement \(\Pi_X\) is
clause~(vi) of the same lemma; clause~(vii) shows no additive linear
\(L:\Gamma_X^{\mathrm{form}}\to\Gamma_{\mathrm{gram}}\) can satisfy
\(B=\delta L\), since \(B(c,c)=2\Pi_X(c)\) is non-zero on charges with
non-trivial Mukai self-pairings.
\end{proof}

\begin{remark}[Cleanly stated polarization identity]
\label{rem:cleanly-stated-polarization-identity}
\ClaimStatusDerivation\
The polarization identity
\(B(c,c)=2\Pi_X(c)\) is the symmetric companion to
\(\Pi_X(c+c')=\Pi_X(c)+\Pi_X(c')+B(c,c')\); evaluated at
\(c=c'\) the latter gives
\(\Pi_X(2c)=2\Pi_X(c)+B(c,c)\), and \(\Pi_X(2c)=4\Pi_X(c)\) by
quadraticity of \(\Pi_X\), hence \(B(c,c)=2\Pi_X(c)\). This is purely
chart-data structure; no Heisenberg or central extension hypothesis on
the source is invoked.
\end{remark}

\subsection{Normal-ordered Gram map}
\label{subsec:normal-ordered-gram-map-appendixC}

\begin{definition}[Normal-ordered charge lattice]
\label{def:normal-ordered-charge-lattice-appendixC}
\ClaimStatusDefinition\
Set
\[
\widehat\Gamma_X
:=\Gamma_X^{\mathrm{form}}\oplus_B\Gamma_{\mathrm{gram}},
\]
the underlying set
\(\Gamma_X^{\mathrm{form}}\times\Gamma_{\mathrm{gram}}\) equipped with
the abelian-group law
\[
(c,T)\star(c',T'):=\bigl(c+c',\,T+T'+B(c,c')\bigr).
\]
The identity is \((0,0)\); the inverse of \((c,T)\) is
\((-c,\,-T+B(c,c))=(-c,\,-T+2\Pi_X(c))\). The split section
\[
s:\Gamma_X^{\mathrm{form}}\longrightarrow\widehat\Gamma_X,
\qquad s(c):=(c,\Pi_X(c)),
\]
is a homomorphism of abelian groups; the raw placement
\[
i_0:\Gamma_X^{\mathrm{form}}\longrightarrow\widehat\Gamma_X,
\qquad i_0(c):=(c,0),
\]
is not a homomorphism unless \(B\equiv0\).
\end{definition}

\begin{proposition}[Normal-ordered Gram homomorphism]
\label{prop:normal-ordered-gram-homomorphism-appendixC}
\ClaimStatusProvedElsewhere\
The map
\[
\overline\Pi_X:\widehat\Gamma_X\longrightarrow\Gamma_{\mathrm{gram}},
\qquad
\overline\Pi_X(c,T):=\Pi_X(c)-T,
\]
is a homomorphism of abelian groups, \(\overline\Pi_X\colon
(\widehat\Gamma_X,\star)\to(\Gamma_{\mathrm{gram}},+)\); equivalently,
\[
\overline\Pi_X\bigl((c,T)\star(c',T')\bigr)
=\overline\Pi_X(c,T)+\overline\Pi_X(c',T').
\]
On the split section,
\(\overline\Pi_X(s(c))=0\); on the raw placement,
\(\overline\Pi_X(i_0(c))=\Pi_X(c)\).
\end{proposition}

\begin{proof}
Direct calculation,
\begin{align*}
\overline\Pi_X((c,T)\star(c',T'))
&=\overline\Pi_X(c+c',T+T'+B(c,c'))\\
&=\Pi_X(c+c')-T-T'-B(c,c')\\
&=\bigl(\Pi_X(c)+\Pi_X(c')+B(c,c')\bigr)-T-T'-B(c,c')\\
&=\bigl(\Pi_X(c)-T\bigr)+\bigl(\Pi_X(c')-T'\bigr)\\
&=\overline\Pi_X(c,T)+\overline\Pi_X(c',T'),
\end{align*}
using Lemma~\ref{lem:gram-additivity-defect-appendixC}. The split-section
and raw-placement identities are the definitions.
\end{proof}

The finite formal packet
\(\texttt{certificates/charge/mukai\_gram\_cocycle}\) records explicit
rank-one Mukai-slice rows for the three simple Gram shadows
\(\delta_1,\delta_2,\delta_3\) and one test charge.  Its verifier checks
the fixed Mukai sign convention, integrality of \(\Pi_X\), symmetry and
bilinearity of \(B\), the identity
\[
B(c,c')=\Pi_X(c+c')-\Pi_X(c)-\Pi_X(c'),
\]
the polarization identity \(B(c,c)=2\Pi_X(c)\), additivity of the split
section \(s(c)=(c,\Pi_X(c))\), non-additivity of the raw placement
\((c,0)\), and the homomorphism property of \(\overline\Pi_X\).  This
packet is formal chart arithmetic; it is not a finite HN charge window,
a compact source construction, a Hall bracket computation, a Pfaffian
orientation, an \(O2\) wall atlas, a mirror discriminant, or a protected
trace.

\begin{proposition}[Compact Mukai representatives for the three simple wall charges]
\label{prop:type-ii-compact-mukai-representatives}
\textnormal{[constructed in the compact \(K3\) Mukai chart]}
Let \(S\) be a smooth projective \(K3\) surface.  Choose a closed point
\(p\in S\) and a zero-dimensional closed subscheme \(Z\subset S\) of
length \(2\).  With the convention
\[
v(E)=\ch(E)\sqrt{\td(S)},\qquad \sqrt{\td(S)}=(1,0,1),
\]
one has
\[
v(\mathcal I_Z)=(1,0,-1),\qquad
v(\mathcal I_p)=(1,0,0),\qquad
v(\mathcal O_p)=(0,0,1).
\]
Define compact Mukai-pair representatives in
\(\widetilde H(S,\Z)\oplus\widetilde H(S,\Z)\) by
\[
x_{\delta_1}:=(\mathcal I_Z,\mathcal I_p),\qquad
x_{\delta_2}:=(\mathcal I_p,\mathcal I_Z),\qquad
x_{\delta_3}:=(\mathcal I_p,\mathcal O_p).
\]
Then
\[
\Pi_X(v(x_{\delta_1}))=(1,1,0),\qquad
\Pi_X(v(x_{\delta_2}))=(0,1,1),\qquad
\Pi_X(v(x_{\delta_3}))=(0,-1,0).
\]
Thus the three triples of
Proposition~\ref{prop:type-ii-wall-images} have compact representatives
at the Mukai-chart level.  This construction does not supply retained
HN semistability, a compact Hall source basis, reduced orientations,
vanishing-cycle summands, \(E\)-quotient data, O2 wall charts, Pfaffian
signs, or protected integration.
\end{proposition}

\begin{proof}
For a zero-dimensional closed subscheme \(W\subset S\) of length \(r\),
\(\ch(\mathcal I_W)=1-r[\mathrm{pt}]\), hence
\[
v(\mathcal I_W)=(1,0,1-r).
\]
For \(r=2\) this gives \(v(\mathcal I_Z)=(1,0,-1)\), and for \(r=1\)
it gives \(v(\mathcal I_p)=(1,0,0)\).  The skyscraper sheaf has
\(\ch(\mathcal O_p)=(0,0,1)\), hence
\(v(\mathcal O_p)=(0,0,1)\).

Use the rank-one Mukai pairing
\[
\bigl((r,d,s),(r',d',s')\bigr)=2dd'-rs'-r's.
\]
The three pair representatives have Mukai vectors
\[
\begin{array}{c|cc}
&Q&P\\
\hline
x_{\delta_1}&(1,0,-1)&(1,0,0)\\
x_{\delta_2}&(1,0,0)&(1,0,-1)\\
x_{\delta_3}&(1,0,0)&(0,0,1).
\end{array}
\]
Therefore
\[
\Pi_X(Q,P)=\Bigl(\frac{(Q,Q)}2,(Q,P),\frac{(P,P)}2\Bigr)
\]
gives respectively
\[
\Bigl(\frac2 2,1,\frac0 2\Bigr)=(1,1,0),\qquad
\Bigl(\frac0 2,1,\frac2 2\Bigr)=(0,1,1),\qquad
\Bigl(\frac0 2,-1,\frac0 2\Bigr)=(0,-1,0).
\]
All three sheaves are coherent sheaves on the projective surface \(S\),
so their supports are proper over \(\C\).  The final sentence records
the extra structure not produced by the Mukai-vector calculation.
\end{proof}

The finite packet
\[
\texttt{certificates/charge/type\_ii\_compact\_mukai\_representatives}
\]
verifies
\(\texttt{TYPE\_II\_COMPACT\_MUKAI\_REPRESENTATIVES\_VERIFIED}\): the
three compact Mukai-pair representatives above have the required
Gram-degree shadows.  Its scalar firewall excludes reading these rows
as compact Hall source basis vectors or retained O2 wall objects.

\begin{remark}[Polarization defect]
\label{rem:resolution-polarization-defect-appendixC}
\ClaimStatusProved\
The composite
\(\overline\Pi_X\circ i_0=\Pi_X\) records the quadratic Gram-degree
shadow of the bare upstairs charge; the composite
\(\overline\Pi_X\circ s\equiv0\) records the split-section convention
that absorbs the entire Gram-degree into the central coordinate
\(T\). The choice of section is the choice of normal ordering: bare
placements collect cross-pairing \(B(c_i,c_j)\) in the central
coordinate (Lemma~\ref{lem:gram-additivity-defect-appendixC}), while
\(s\)-placements record \(\overline\Pi_X=0\) on every component and the
Gram-degree information lives entirely in the upstairs \(\Pi_X(c)\). The
two presentations agree on the homomorphism \(\overline\Pi_X\). This is
the additive group law on which a hypothetical compact Hall product can
descend without a chain-level associator at the lattice level.
\end{remark}

\begin{remark}[Chart data, not Hall support]
\label{rem:chart-data-not-hall-support-appendixC}
\ClaimStatusProved\
The lattice \(\widehat\Gamma_X\) is the formal additive group on which
charges and their Gram-degree shadows are tracked. It is not the
algebraic effective semigroup \(\Gamma_X^{\mathrm{eff}}\) of curve
classes representable by closed subschemes \(Y\subset X\); membership
in the effective sub-semigroup is a separate algebraic-effective
condition (Definition~\ref{def:algebraic-effective-hall-support}). The
Beilinson chart layer \(\alpha\) (charge data) and \(\gamma\) (Mukai
lattice) license \(\widehat\Gamma_X\) as ambient; the Hall-support
restriction is part of the categorical-input layer at higher Beilinson
level, not the Mukai--Gram chart datum.
\end{remark}

\subsection{Three Igusa names and the Mukai--Gram chart}
\label{subsec:three-igusa-names-chart-data}

Three Igusa names are distinct mathematical objects and remain distinct under
the Mukai--Gram chart datum:
\begin{itemize}
\item the automorphic section \(\mathcal D_X=\Delta_5\), a section of
  the weight-\(5\) automorphic line over the genus-\(2\) Siegel space
  with Maass character \(\nu_{\Delta_5}\);
\item the BKM denominator
  \(\mathrm{den}(\mathfrak g_{\Delta_5})=64^{-1}\Delta_5(2Z)\)
  on the root lattice \(\Lambda_{II}^{2,1}\); and
\item the compact protected Pfaffian
  \(\operatorname{Pf}_{\mathrm{prot}}(\mathfrak D_X)\), a section of
  \(\mathcal L^5\otimes\nu_{\Delta_5}\) on the conditional pro-object
  \(\mathfrak D_X^{\mathrm{DI}}\) of \S5.
\end{itemize}
The Mukai--Gram dictionary supplies the Gram-degree shadow
\(\Pi_X:\Gamma_X^{\mathrm{form}}\to\Gamma_{\mathrm{gram}}\) on which
the cusp-Fourier expansion of \(\Delta_5\), the root-degree map of
\(\mathrm{den}(\mathfrak g_{\Delta_5})\), and the wall data of
\(\operatorname{Pf}_{\mathrm{prot}}\) all read.
The Mukai pairing on \(\widetilde H(S,\Z)\) is symmetric in both
algebraic and transcendental parts; the ungraded Mukai-form classical
limit of an associated Yangian-type construction is the orthogonal Lie
algebra \(\mathfrak{so}(4,20)\), not the orthosymplectic
\(\mathfrak{osp}(4\mid20)\). The Hodge \(\Z/2\)-super-extension that
appears in Yangian conventions imposes a separate
\(\Z/2\)-grading on the algebraic side; it does not enter the lattice
pairing above. The Bryan/Oberdieck--Pandharipande/
Maulik--Toda conventions \cite{OPand,BRYAN,MaulikTodaGV} take \(\beta\)
algebraic and \(d\) elliptic-degree, agreeing with
Theorem~\ref{thm:d6-d2-d0-dictionary-appendixC}.

The Mukai--Gram chart datum does not touch the orientation character
\(\epsilon_o\) of
the Weyl-transport condition \ref{def:O1-plus-weyl-transport}, the Pfaffian
section \(\operatorname{Pf}_{\mathrm{prot}}(\mathfrak D_X)\), or the
recognition theorem of
Theorem~\ref{thm:pfaffian-dirac-finite-stage}; those structures live
at higher Beilinson levels and depend on the chart datum.
